%% ============================================================
%%  Adaptive Vision-Language-Action Systems for Task Changeover
%%  in Collaborative Manufacturing: A Review
%%
%%  Compile:  pdflatex main  ->  bibtex main (not needed, bib is
%%            embedded)  ->  pdflatex main  ->  pdflatex main
%% ============================================================

\documentclass[conference]{IEEEtran}
\IEEEoverridecommandlockouts

\usepackage{cite}
\usepackage{amsmath,amssymb}
\usepackage{graphicx}
\usepackage{booktabs}
\usepackage{array}
\usepackage{multirow}
\usepackage{url}
\usepackage[hidelinks]{hyperref}
\usepackage{textcomp}

% Column type for wrapped, top-aligned table cells
\newcolumntype{L}[1]{>{\raggedright\arraybackslash}p{#1}}

\begin{document}

\title{Adaptive Vision-Language-Action Systems for Task\\
Changeover in Collaborative Manufacturing: A Review\\
and a Proposed Architecture}

\author{\IEEEauthorblockN{Batch 03}
\IEEEauthorblockA{\textit{Department of Computer Science and Engineering (AI and ML)}\\
\textit{Vardhaman College of Engineering}\\
Hyderabad, India}
\and
\IEEEauthorblockN{Dr.~Ramachandro Majji}
\IEEEauthorblockA{\textit{Associate Professor}\\
\textit{Department of Computer Science and Engineering (AI and ML)}\\
\textit{Vardhaman College of Engineering}\\
Hyderabad, India}
}

\maketitle

% ============================================================
\begin{abstract}
Vision-language-action (VLA) models let a robot map a camera image and a
natural language instruction directly to motor commands. Recent open
models such as OpenVLA reach 70.6\% mean success across 29 manipulation
tasks while remaining fine-tunable on a single GPU. All of this work
shares one assumption: the instruction holds constant for the length of
the episode. On a high-mix, low-volume assembly line that assumption
breaks several times a shift, because the operator changes the product
variant while the robot is partway through building the previous one.

This paper reviews vision-language robotics for collaborative
manufacturing through the specific lens of task change. We group twenty
representative works into generalist VLA policies, language grounding
and ambiguity resolution, LLM-based planning and replanning, and
parameter-efficient continual adaptation, and we show that each group
leaves the same thing untouched. Ambiguity work resolves which object an
instruction refers to while holding the goal fixed. Replanning work
reacts to execution failure while holding the goal fixed. Recent task-switching and undo-on-correction work handles the case where the goal itself is withdrawn and replaced mid-episode at the policy or plan level, without reasoning about which completed steps can stay, which must come off, and when a change must be refused because a step cannot be undone.

We then propose TaskFlux, an adaptation layer that sits above an
existing VLA rather than replacing it. Its four components are a
changeover intent classifier, a partial-assembly state reconciler that
decides which completed steps to keep, undo or discard, a feasibility
gate that can refuse a request and explain why, and a LoRA-based
specialised policy head whose adapters merge into the base model only
after passing a retention check. Because we found no benchmark measuring rework or refusal for
mid-episode task change, we also define an evaluation protocol with five
metrics.
\end{abstract}

\begin{IEEEkeywords}
vision-language-action models, human-robot collaboration, task
replanning, high-mix low-volume manufacturing, parameter-efficient
adaptation
\end{IEEEkeywords}

% ============================================================
\section{Introduction}

High-mix, low-volume manufacturing is the awkward middle of production.
Volumes are too low to justify a dedicated line and too high to do by
hand. Medical device assembly, custom electronics and aerospace
subassembly all live there, and in all three the product variant changes
often, sometimes several times in one shift.

Industrial robots handle this badly. A variant switch means a stop, a
reprogram, a revalidation and a restart, and the cost of that downtime is
what keeps a lot of this work manual. Collaborative robots were supposed
to help. They do help with safety and with physical proximity to the
operator, but the reprogramming problem is untouched. A cobot that needs
a technician with a teach pendant every time the part changes is not
flexible in any way that matters to a plant manager.

Vision-language-action models looked like the answer. Train one policy on
enough demonstrations, condition it on a natural language instruction,
and it should generalise to new objects and new phrasings without
retraining. RT-2 showed the idea worked \cite{rt2}. OpenVLA showed it
could be open, reproducible, and cheap enough to adapt on a single A100
in ten to fifteen hours using LoRA at rank 32, which trains 1.4\% of the
model's weights and still matches full fine-tuning \cite{openvla}.

But look at how these systems are actually evaluated. An episode begins,
an instruction is provided, the policy runs to completion or failure, the
episode ends. The instruction is an input, supplied once. There is no
path in the architecture for it to change while the robot is working.

That is exactly what happens in a real changeover. And the moment it
happens, a second problem appears that the VLA work we reviewed does not treat: the
workspace is no longer empty. Parts have already been placed. Screws are
already torqued. A plan generated fresh from the current camera frame
will not account for any of that, and executing it will produce either a
collision or a hybrid assembly that matches neither specification.

This paper does two things. First it reviews the field with that specific
failure in mind, to establish that the gap is real rather than assumed.
Second it proposes an architecture for closing it, along with the
evaluation protocol that would be needed to test any such system.

\subsection{Contributions}

\begin{enumerate}
\item A structured review of twenty works across four themes, each
assessed against a single question: what happens if the task changes
mid-episode.
\item A three-way distinction between referential ambiguity, execution
failure and task changeover, which have been conflated in the literature
and which require different mechanisms.
\item TaskFlux, an adaptation layer over an existing VLA, whose central
component is a partial-assembly state reconciler.
\item A five-metric evaluation protocol for mid-episode task change,
since no current benchmark measures it.
\end{enumerate}

Section~\ref{sec:background} covers background. Section~\ref{sec:method}
states the review method. Section~\ref{sec:review} is the thematic
review. Section~\ref{sec:gap} is the gap analysis.
Section~\ref{sec:arch} presents TaskFlux. Section~\ref{sec:eval} gives
the evaluation protocol. Section~\ref{sec:open} lists open problems.

% ============================================================
\section{Background}
\label{sec:background}

\subsection{From vision-language models to vision-language-action models}

CLIP established that contrastive pretraining on image-text pairs
produces representations that transfer to tasks the model was never
trained on \cite{clip}. Robotics picked this up quickly. CLIPort fused
CLIP semantics with a spatial stream for language-conditioned pick and
place \cite{cliport}, and PerAct extended the idea into voxel space
\cite{peract}.

The step to a VLA came when RT-2 treated actions as text \cite{rt2}.
Discretise each dimension of the robot action into bins, assign each bin
a token, and the whole control problem becomes next-token prediction over
a vocabulary that happens to include motor commands. That means every
technique developed for training large language models applies directly
to training robot policies.

OpenVLA is the open realisation of that idea and the backbone this
project uses \cite{openvla}. It is a 7B model built on Prismatic-7B,
which pairs Llama~2 with a two-part visual encoder that concatenates
SigLIP and DINOv2 features. The authors chose that pairing deliberately:
SigLIP carries high-level semantics, DINOv2 carries low-level spatial
detail, and spatial detail is what a manipulation policy needs. Each of
the seven action dimensions is discretised into 256 bins set by the 1st
and 99th quantiles of the training distribution, and those 256 tokens
overwrite the 256 least-used entries in the Llama tokenizer.

Trained on 970k episodes from Open X-Embodiment \cite{openx}, it reaches
70.6\% $\pm$ 3.2\% mean success across 29 tasks, which is 16.5 absolute
points above RT-2-X despite having seven times fewer parameters.

\subsection{What the OpenVLA authors say is missing}

Their own limitations section is worth taking seriously, because it
constrains what any system built on top can do.

The model sees one image. No proprioception, no wrist camera, no
observation history. Inference throughput is too low for high-frequency
control; the authors note it cannot support something like ALOHA at
50~Hz. And on narrow dexterous tasks, Diffusion Policy still produces
smoother trajectories, which the authors attribute to action chunking and
temporal smoothing that OpenVLA does not do.

For our purposes the throughput limit is the binding one. A changeover
that takes a few seconds to reason about is acceptable when changeovers
happen a handful of times per shift. It would not be acceptable for
anything reactive.

\subsection{Parameter-efficient adaptation}

LoRA freezes the base weights and learns a low-rank update on each linear
layer \cite{lora}. On OpenVLA the numbers are unusually favourable. Rank
32 gives 68.2\% against 69.7\% for full fine-tuning, trains 97.6M
parameters instead of 7.19B, and fits in 59.7~GB at batch size 16 where
full fine-tuning needs 163.3~GB sharded across two GPUs. Rank had almost
no effect on the result, so 32 is a reasonable default rather than a
tuned choice.

That is what makes a per-variant adapter practical. It is also what makes
the merge-back question interesting, because if adapters are cheap you
will accumulate a lot of them and you need a policy for what to do with
the pile.

% ============================================================
\section{Review Method}
\label{sec:method}

We started from six sources compiled during the project's initial survey
and expanded outward along citation edges from two anchors: OpenVLA as
the technical backbone \cite{openvla} and the survey of Fan et al. in
Frontiers of Engineering Management as the domain reference
\cite{fansurvey}. That survey itself screened Web of Science, Scopus and
IEEE Xplore for combinations of ``human-robot'' and ``vision language''
over 2020 to 2024 and retained 109 papers, so it provides reasonable
coverage of the manufacturing side without us repeating the search.

Inclusion required a paper to do at least one of the following: propose
or evaluate a policy that consumes both visual and language input,
propose a planning or replanning mechanism driven by language, or propose
an adaptation method we intend to use. Twenty works met that bar and
appear in Table~\ref{tab:review}.

\subsection{One source excluded}

One source was excluded after screening. A 2025 article in the Journal of
Computer Science and Software Applications carries the title
``Vision-Language Models for Human-Robot Collaboration: Real-Time Task
Understanding and Execution'' over an abstract and keyword list that
describe CNN-LSTM stock price prediction on the CSI 300 index. The body
then discusses robots. The mismatch is present in the published version
of record, which indicates the manuscript was never reviewed.

Its reported figures are round numbers with no variance, from twenty
trials per condition, and we do not treat them as evidence. We mention
the exclusion rather than quietly dropping the paper, because venues of
this kind now appear routinely in search results for this topic and other
groups will encounter the same article.

% ============================================================
\section{Thematic Review}
\label{sec:review}

\subsection{Generalist VLA policies}

RT-1 showed one transformer could absorb hundreds of manipulation tasks
\cite{rt1}. RT-2 added web-scale semantic knowledge by co-fine-tuning a
pretrained VLM on robot trajectories \cite{rt2}. Open X-Embodiment
aggregated over sixty datasets across twenty-two embodiments and
demonstrated positive transfer between them \cite{openx}. OpenVLA made
the recipe open \cite{openvla}. Octo showed the same territory can be
reached with a diffusion action head and modular input tokenisers
\cite{octo}.

QUAR-VLA is useful as a control on the embodiment question. Ding et al.
built QUART for quadrupeds over an 11-dimensional command space and
trained on QUARD, 259k simulated plus 3k real episodes \cite{quarvla}.
The gap between QUART and the CLIP, R3M and VC-1 baselines widens as
tasks get harder, and on the crawl and unload tasks the baselines score
exactly zero while QUART reaches 0.32 and 0.12. The formulation clearly
transfers off arms.

What all six share is this: the instruction is bound once, at the start
of the episode, and the architecture provides no way to rebind it.

\subsection{Language grounding and ambiguity in HRC}

Fan and Zheng \cite{fanzheng} is the closest published work to this
project and should be the primary baseline. They target under-specified
operator instructions in collaborative assembly, the ``hand me that one''
case, and use vision-language guidance to pick the intended referent.

It is worth being precise about what that solves and what it does not.
The operator and the robot already agree on what is being built. The
uncertainty is over which physical object satisfies a step of an agreed
plan. Resolve it and execution continues along the same plan. This is
genuinely useful and genuinely different from what happens when the
operator says ``we are building the B variant now'', where the plan
itself is no longer valid.

OWG approaches grounding from the grasp side. Tziafas and Kasaei prompt
GPT-4V with visual markers over segmented objects and candidate grasps,
in three zero-shot stages: referring segmentation, grounded grasp
planning, then grasp ranking by contact reasoning \cite{owg}. On real
hardware with two UR5e arms they reach 83.3\% on isolated unseen objects
and 50.0\% in clutter, beating both a supervised baseline and an LLM
planner baseline in every cell.

Their limitations section carries a lesson we adopted. Being modular, OWG
``suffers from error cascading effects introduced by the segmentor and
grasp synthesis models''. A segmentation error becomes a grasp error
becomes a task failure, and there is no path to recover. That is a direct
argument for keeping the low-level policy monolithic and putting the
language reasoning above it rather than in series with perception.

\subsection{LLM planning and replanning}

SayCan pairs an LLM's estimate of what a skill contributes toward a goal
with an affordance model's estimate of whether that skill will succeed
here and now, and multiplies them \cite{saycan}. The feasibility gate is
the part we generalise: our safety module does the same job but checks
process constraints, not just affordances.

Code as Policies has the LLM emit executable code that calls perception
and control APIs \cite{cap}, which makes the plan inspectable and
verifiable instead of free text. That is a property we want for a plan
that a human operator may need to approve.

Inner Monologue is the canonical replanning system \cite{innermono}.
Success detectors, scene descriptors and human feedback are all injected
back into the LLM's context as text, and the plan is revised when reality
diverges from prediction.

Read those three together and the pattern is clear. The goal is a
constant. Planning generates a route to it, replanning generates a new
route when the old one fails, and no mechanism exists for the goal to be
replaced. That is the structural gap.

\subsection{Parameter-efficient adaptation and continual learning}

LoRA \cite{lora} and QLoRA \cite{qlora} make per-variant adaptation cheap
enough to do on-premise. Model soups \cite{soups} and task arithmetic
\cite{taskarith} show that merging in weight space works and can even be
reversed, which matters if an adapter turns out to have hurt something.
Elastic weight consolidation gives the theoretical account of why
unconstrained sequential fine-tuning destroys earlier competence
\cite{ewc}.

None of this has been assembled into a loop that runs on a factory floor
with a validation gate in the middle. That assembly is part of what we
are proposing.

\subsection{The manufacturing view}

The survey of Fan et al. \cite{fansurvey} is where the domain
requirements are stated most clearly, and two of its future-work sections
read like a specification for this project.

Section 6.2 says VLM-based task planning is focused on static scenes and
that ``real-time task planning in dynamic scenes remains an unresolved
issue''.

Section 6.6, titled ``Dynamic task adaptation and unsupervised
evaluation'', says current methods ``always rely on an assumption of an
ideal training environment'', that real deployment ``necessitates human
operators for continuous monitoring and intervention'', and that this
``significantly hinders real applications''. It then argues robots ``must
evolve to autonomously adapt to dynamic environments and tasks'' and
calls for ``a continuous learning mechanism'' allowing them ``to adapt
and improve their performance autonomously over time based on new
experiences and feedback''.

That is our problem statement, written by someone else, in a survey of
109 papers, with no solution attached.

% ============================================================
% Wide comparative table
% ============================================================
\begin{table*}[!t]
\caption{Comparative review of vision-language robotics for
collaborative manufacturing, assessed against mid-episode task change}
\label{tab:review}
\centering
\footnotesize
\renewcommand{\arraystretch}{1.15}
% Widths sum to 15.35cm + 5 inter-column gaps (2\tabcolsep each,
% ~2.1cm total) = ~17.45cm, inside the 18.1cm IEEE double-column
% textwidth. Adjust these six numbers together if you edit the table.
\begin{tabular}{@{}L{0.45cm} L{2.4cm} L{0.75cm} L{3.85cm} L{3.85cm} L{4.05cm}@{}}
\toprule
\textbf{\#} & \textbf{Work} & \textbf{Year} & \textbf{Approach} &
\textbf{Reported result} & \textbf{Gap w.r.t. task changeover} \\
\midrule
\multicolumn{6}{@{}l@{}}{\textit{A. Generalist vision-language-action policies}}\\
\midrule
1 & OpenVLA \cite{openvla} & 2024 &
7B VLA. Llama~2 with fused SigLIP and DINOv2 encoder; actions
discretised into 256 bins over unused tokens &
70.6\% $\pm$ 3.2\% over 29 tasks; 16.5 pts above RT-2-X at 7$\times$
fewer parameters; LoRA $r{=}32$ gives 68.2\% training 1.4\% of weights &
Instruction fixed for the episode. Single image, no proprioception, no
history. Control rate too low for reactive use \\
2 & QUAR-VLA \cite{quarvla} & 2024 &
VLA for quadruped locomotion and manipulation over an 11-dimensional
command space &
QUART 0.66 on seen tasks vs 0.44 CLIP, 0.46 VC-1; baselines score 0.0 on
crawl and unload where QUART reaches 0.32 and 0.12 &
Task set enumerated in advance. No mechanism for a task that changes
during execution \\
3 & RT-2 \cite{rt2} & 2023 &
Web-pretrained VLM co-fine-tuned on robot trajectories, actions emitted
as text tokens &
Strong emergent semantic generalisation reported &
Closed weights. Fixed-task assumption. No replanning layer \\
4 & RT-1 \cite{rt1} & 2022 &
Transformer policy over image and instruction tokens; 130k
demonstrations, 700 tasks &
High multi-task success on a mobile manipulator &
Task identity supplied at episode start and never revised \\
5 & Open X-Embodiment \cite{openx} & 2024 &
Cross-embodiment dataset aggregation and joint policy training &
Positive transfer across 22 embodiments &
Dataset contribution. No task adaptation mechanism \\
6 & Octo \cite{octo} & 2024 &
Transformer policy with diffusion action head and modular tokenisers &
Competitive with generalist baselines; flexible fine-tuning &
Same fixed-goal assumption \\
\midrule
\multicolumn{6}{@{}l@{}}{\textit{B. Language grounding and ambiguity in human-robot collaboration}}\\
\midrule
7 & Fan and Zheng \cite{fanzheng} & 2024 &
Vision-language guided action planning for ambiguity mitigation in
collaborative assembly &
High task success on ambiguous referring instructions &
Resolves \emph{which object} an instruction denotes. The goal state
itself is fixed throughout. Closest prior work and the correct baseline \\
8 & OWG \cite{owg} & 2024 &
GPT-4V with Mask-RCNN and GR-ConvNet; referring segmentation, grounded
grasp planning, contact-reasoning grasp ranking, all zero-shot &
Sim seen/unseen isolated 78.0/82.0, cluttered 62.0/66.0; real isolated
83.3/66.6, cluttered 50.0/50.0 &
Single-object grasping, not multi-step assembly. Authors flag error
cascading between modules \\
9 & CLIPort \cite{cliport} & 2021 &
Two-stream policy fusing CLIP semantics with a spatial Transporter
stream &
Strong sample efficiency on language-conditioned pick and place &
Tabletop rearrangement only, one instruction per episode \\
10 & PerAct \cite{peract} & 2022 &
Perceiver-based voxel policy for language-conditioned manipulation &
Multi-task success from few demonstrations &
No plan-level reasoning, no adaptation to a changed goal \\
11 & CLIP \cite{clip} & 2021 &
Contrastive image-text pretraining on 400M pairs &
Strong zero-shot transfer &
Not a robotics model. Provides similarity, not action \\
\midrule
\multicolumn{6}{@{}l@{}}{\textit{C. LLM task planning and replanning}}\\
\midrule
12 & SayCan \cite{saycan} & 2022 &
LLM scores candidate skills, affordance model scores feasibility, product
selects the next skill &
Long-horizon instruction following on a real mobile manipulator &
Skill library fixed. Plan generated once from a goal that never changes \\
13 & Code as Policies \cite{cap} & 2023 &
LLM writes executable policy code calling perception and control APIs &
Generalises to novel instructions through code composition &
Generates a program for a stated goal. No notion of revising it
mid-execution \\
14 & Inner Monologue \cite{innermono} & 2022 &
Closed-loop replanning from textual feedback: success detectors, scene
descriptors, human input &
Improved long-horizon success from feedback grounding &
Replans when execution diverges from prediction. The goal is never
withdrawn and replaced \\
15 & Fan et al. survey \cite{fansurvey} & 2024 &
Systematic survey; Web of Science, Scopus, IEEE Xplore; 2020--2024 &
109 papers; taxonomy across planning, navigation, manipulation, skill
transfer &
Survey only. Section 6.6 names this gap and calls for a continuous
learning mechanism, but builds nothing \\
\midrule
\multicolumn{6}{@{}l@{}}{\textit{D. Parameter-efficient adaptation and continual learning}}\\
\midrule
16 & LoRA \cite{lora} & 2022 &
Low-rank update matrices injected into frozen linear layers &
Matches full fine-tuning at a fraction of trainable parameters &
Designed for offline adaptation to a static target task \\
17 & QLoRA \cite{qlora} & 2023 &
LoRA over a 4-bit quantised base model &
Large memory reduction at near-parity quality &
Same offline assumption \\
18 & Model soups \cite{soups} & 2022 &
Averaging weights of independently fine-tuned models &
Averaging can beat the best single member &
Merging is unconditional. No gating on retention or safety \\
19 & Task arithmetic \cite{taskarith} & 2023 &
Add and subtract task vectors in weight space to edit behaviour &
Behaviour can be composed and removed arithmetically &
Not validated in a control or safety-critical setting \\
20 & EWC \cite{ewc} & 2017 &
Penalise change to parameters important for earlier tasks &
Reduces catastrophic forgetting &
Predates VLAs. Never tested on a manipulation policy \\
\bottomrule
\end{tabular}
\end{table*}

% ============================================================
\section{Gap Analysis}
\label{sec:gap}

\subsection{Three problems that get conflated}

The literature contains two named problems and one unnamed one, set out
in Table~\ref{tab:three}.

\begin{table}[!t]
\caption{Three problems that are routinely conflated}
\label{tab:three}
\centering
\footnotesize
\renewcommand{\arraystretch}{1.25}
\begin{tabular}{@{}L{2.0cm} L{2.6cm} L{1.5cm} L{1.3cm}@{}}
\toprule
\textbf{Problem} & \textbf{Trigger} & \textbf{Goal state} &
\textbf{Addressed by} \\
\midrule
Referential ambiguity & Referent is unclear & Fixed &
\cite{fanzheng} \\
Execution failure & Execution diverged from prediction & Fixed &
\cite{innermono} \\
Task changeover & Operator changed the requirement & Replaced
mid-episode & Policy or plan level only \\
\bottomrule
\end{tabular}
\end{table}

The three need different machinery. Ambiguity needs better grounding.
Failure needs monitoring and recovery. Changeover needs goal-state
comparison and a way to deal with work already done.

\subsection{The partial-state problem}

This is the part that has no precedent at all.

Suppose a plan has $n$ steps and the goal is replaced at step $k$. The
naive response is to regenerate a plan from the new goal and the current
image. That fails, because the current image shows a partially assembled
product and the planner will either treat the placed parts as obstacles
to avoid or ignore them entirely. Either way the output is wrong.

What is needed instead is a comparison between the effects of steps $1$
through $k$ and the requirements of the new goal, producing a partition
into steps to keep, steps to undo, and steps to discard, plus a valid
ordering for the undos. Undo has to respect dependencies, since an outer
cover comes off before the inner bracket, and it has to respect
irreversibility, because cured adhesive and set rivets do not come back
out.

We found nothing in the reviewed literature that does this with a completeness guarantee, although plan repair, selective disassembly planning and undo-on-correction work are adjacent. Rearrangement planning in
classical robotics is adjacent but assumes a symbolic world model with
clean pre- and post-conditions, which is exactly what a VLA does not give
you.

\subsection{Adaptation without forgetting}

A system that learns variant B by fine-tuning on B and then cannot build
variant A any more has not adapted. It has moved. Every project in this
space needs a retention answer and most do not state one.

\subsection{There is no benchmark for rework or refusal}

Every dataset in Table~\ref{tab:review} supplies one instruction per
episode. There is no public benchmark with a mid-episode instruction
change, no metric for how fast a system adapts, and no metric for how
much finished work an adaptation destroyed. Defining these is a
prerequisite for the field, not an afterthought.

% ============================================================
\section{Proposed Architecture: TaskFlux}
\label{sec:arch}

TaskFlux is an adaptation layer over an existing VLA, in six stages.

\subsection{Stage 1: instruction intake}

Operator speech is transcribed and passed to the orchestrator along with
the current camera frame and the execution state, meaning which steps of
the active plan have completed.

\subsection{Stage 2: agentic orchestrator}

An LLM performs four jobs in sequence.

\emph{Intent understanding} produces a structured reading of the
utterance: what the operator wants, which product variant is implied, and
which constraints they stated.

\emph{Changeover classification} assigns the utterance to one of four
types: clarification, parameter edit, changeover, or abort. Only a
changeover triggers reconciliation. This classification is where a large
share of the system's errors will live, and the two error directions have
very different costs, so the classifier should be biased toward asking
rather than guessing.

\emph{Feasibility and safety analysis} checks reachability, collision
against the current workspace occupancy, tool and fixture availability,
and process constraints such as torque sequence and cure time.

The \emph{critique step} reports what the system cannot do and why, in a
sentence the operator can act on, rather than failing silently or
attempting a partial adaptation.

\subsection{Stage 3: state reconciliation}

The novel component. Inputs are the completed prefix of the old plan, the
goal state implied by the new variant, and a process description that
marks each step's reversibility and its dependencies.

Output is a keep set, an undo set ordered in reverse dependency order,
and a discard set, followed by the residual steps of the new plan. If any
required undo touches an irreversible step, the reconciler stops and
escalates to the operator with the specific blocking step named.

\subsection{Stage 4: policy selection}

A known variant, meaning one present in the skill library, routes to the
base OpenVLA policy. An unknown variant instantiates a LoRA adapter at
rank 32, following the OpenVLA authors' own default, and routes to that.

\subsection{Stage 5: execution and validation}

The selected policy executes. A validation module checks each completed
step against its expected post-condition using the perception module, and
checks the finished assembly against the new specification. Novel-variant
executions run under stricter validation and with a lower intervention
threshold.

\subsection{Stage 6: knowledge update}

Successful adaptations are logged to a dynamic knowledge base holding the
skill library, a vector store of past instructions and their resolved
plans, and execution traces.

The merge engine folds an adapter into the base weights only when two
gates pass. The quality gate requires success above threshold across a
minimum number of executions. The retention gate re-evaluates a held-out
set of previously mastered variants after a candidate merge and rejects
the merge if any regresses beyond tolerance. Because merging is done in
weight space \cite{soups,taskarith}, a rejected merge can be backed out.

\subsection{Why the layer sits above the VLA rather than inside it}

Two reasons, both taken from the reviewed papers.

The error-cascading result in OWG \cite{owg} argues against putting more
modules in series with perception and control. Keeping OpenVLA monolithic
in the execution path limits the depth of the cascade.

Practically, the base model stays a stock checkpoint. That means it can
be swapped for a newer VLA without redesigning the adaptation layer,
which matters in a field where the state of the art moves every few
months.

% ============================================================
\section{Evaluation Protocol}
\label{sec:eval}

The core experimental unit is a \emph{changeover episode}. Begin
executing variant A. At a controlled step $k$, issue an instruction.
Record the classification decision, the reconciliation output, any
refusal, the actions executed, and the final state of the assembly.
Holding $k$ fixed per condition is what makes episodes comparable.

\subsection{Metrics}

\textbf{Changeover Success Rate.} Fraction of episodes where the final
assembly matches the new specification, judged against the acceptance
criteria a human inspector would use. Report per condition, not pooled.

\textbf{Adaptation Latency.} Seconds from the end of the utterance to the
first executed action that is correct under the new plan. Break it into
transcription and classification, reconciliation, and policy load. A
total of six seconds means something very different if five of them are
adapter loading than if five are LLM reasoning.

\textbf{Rework Cost.} Undo actions divided by steps completed when the
change arrived. A full restart scores 1.0 by definition. Report
unnecessary undos separately, since that is the reconciler's precision.

\textbf{False Changeover Rate.} Clarifications and parameter edits
misclassified as changeovers. Report missed changeovers separately. The
two are not symmetric: a false changeover destroys completed work, a
missed changeover produces a finished product that does not match the
order. On a real line the second is worse, so the classifier should be
tuned toward over-triggering and the paper should say so rather than
reporting a single accuracy figure that hides the tradeoff.

\textbf{Retention Score.} Success on variant A after the system has
learned variant B and merged the adapter. Without this number, any claim
about continual learning is unsupported.

\subsection{Conditions}

Cross these deliberately rather than sampling at random: when the change
arrives, early against late; what the change requires, being truncation
only, undo required, or irreversible conflict; whether the variant is
known or novel; and how the instruction is phrased, including indirect
phrasings such as ``we are not doing gaskets on this batch'' rather than
only imperatives.

The irreversible-conflict condition is a refusal test. Success there
means the system refused and explained correctly, not that it completed
the task.

\subsection{Baselines}

\textbf{B1.} Stock OpenVLA with the original instruction. No adaptation.
Should fail every changeover episode. This is the floor.

\textbf{B2.} Stock OpenVLA restarted with the new instruction. Clear the
workspace, reset, run from scratch. This is what a factory does today,
and it quantifies the cost of current practice.

\textbf{B3.} LLM replanner with no reconciler. Regenerate from the new
goal and the current frame, execute directly.

\textbf{B4.} Full TaskFlux.

B3 against B4 is the ablation that matters most. If B3 performs as well
as B4, the reconciler is not earning its place and the central claim is
weak. Run it and report it honestly either way. A negative result that is
properly measured is publishable; an unmeasured claim is not.

\subsection{Sample size}

A single-arm cell will not reach the trial counts in the papers we cite.
OWG ran 50 simulated trials per scenario \cite{owg}; OpenVLA ran 500 per
simulation suite \cite{openvla}. We plan for 10 to 15 episodes per
condition on hardware. Report per-condition counts, give raw numbers
alongside percentages, and include standard error. With $n{=}10$ one
extra failure moves the percentage by ten points, and a reader needs to
see that. Do not pool conditions to get a bigger-looking $n$, because the
conditions are the finding.

% ============================================================
\section{Open Problems}
\label{sec:open}

\emph{Reversibility is hand-authored.} The process description that tells
the reconciler which steps can be undone is written by a human. Learning
it from demonstration or from physical interaction is open and hard.

\emph{Undo is harder than assembly.} Extracting a seated connector is
contact-rich, and the OpenVLA authors concede Diffusion Policy is
smoother on that class of task \cite{openvla}. Undo success may end up
being the binding constraint on the whole system.

\emph{Latency compounds.} Transcription, an LLM call, a reconciliation
pass and a policy switch stack in front of a policy already running at
single-digit hertz.

\emph{Trust and authority are unresolved.} If the robot refuses an
operator instruction, who overrides whom, and what gets logged. A
technical paper can raise this but not settle it.

\emph{Weight-space merging is poorly understood for control policies.}
Model soups and task arithmetic were validated on classification, not on
a policy where a small behavioural regression can mean a collision.

% ============================================================
\section{Conclusion}

The vision-language-action literature has converged on a formulation that
works: condition a pretrained vision-language model on robot
demonstrations, emit actions as tokens, and get generalisation across
objects and phrasings that hand-programmed systems never had. OpenVLA
made that formulation open and cheap to adapt.

The formulation carries one assumption that most of it leaves implicit, which
is that the task is fixed for the duration of the episode. In high-mix
manufacturing it is not, and the failure that follows is not a perception
failure or a grounding failure. It is a system that keeps confidently
building the wrong product.

We reviewed twenty works against that question and found the gap
consistent across all four themes. We separated task changeover from the
two problems it is usually confused with, and argued that it needs
different machinery, specifically a way to reconcile a half-built
workspace against a replaced goal.

TaskFlux is our proposal for that machinery, and the protocol in
Section~\ref{sec:eval} is our proposal for how anyone would know whether
it works. Neither is validated yet. The next phase of this project builds
the reconciler and runs the changeover episodes described above on a
single-arm assembly cell.

% ============================================================
\begin{thebibliography}{20}

\bibitem{openvla}
M.~J. Kim, K.~Pertsch, S.~Karamcheti, \emph{et al.},
``OpenVLA: An open-source vision-language-action model,''
in \emph{Proc. 8th Conf. Robot Learning (CoRL)}, 2024.

\bibitem{rt2}
A.~Brohan \emph{et al.},
``RT-2: Vision-language-action models transfer web knowledge to robotic
control,''
in \emph{Proc. 7th Conf. Robot Learning (CoRL)}, 2023, pp. 2165--2183.

\bibitem{rt1}
A.~Brohan \emph{et al.},
``RT-1: Robotics transformer for real-world control at scale,''
arXiv preprint, 2022.

\bibitem{openx}
A.~O'Neill \emph{et al.},
``Open X-Embodiment: Robotic learning datasets and RT-X models,''
in \emph{Proc. IEEE Int. Conf. Robotics and Automation (ICRA)},
Yokohama, Japan, 2024, pp. 6892--6903.

\bibitem{quarvla}
P.~Ding, H.~Zhao, W.~Song, \emph{et al.},
``QUAR-VLA: Vision-language-action model for quadruped robots,''
in \emph{Proc. European Conf. Computer Vision (ECCV)}, 2024.

\bibitem{owg}
G.~Tziafas and H.~Kasaei,
``Towards open-world grasping with large vision-language models,''
in \emph{Proc. Conf. Robot Learning (CoRL)}, 2024.

\bibitem{fansurvey}
J.~Fan, Y.~Yin, T.~Wang, W.~Dong, P.~Zheng, and L.~Wang,
``Vision-language model-based human-robot collaboration for smart
manufacturing: A state-of-the-art survey,''
\emph{Frontiers of Engineering Management}, 2024.

\bibitem{fanzheng}
J.~Fan and P.~Zheng,
``A vision-language-guided robotic action planning approach for ambiguity
mitigation in human-robot collaborative manufacturing,''
\emph{Journal of Manufacturing Systems}, vol.~74, pp.~1009--1018, 2024.

\bibitem{clip}
A.~Radford \emph{et al.},
``Learning transferable visual models from natural language
supervision,''
in \emph{Proc. Int. Conf. Machine Learning (ICML)}, 2021.

\bibitem{cliport}
M.~Shridhar, L.~Manuelli, and D.~Fox,
``CLIPort: What and where pathways for robotic manipulation,''
in \emph{Proc. Conf. Robot Learning (CoRL)}, 2021.

\bibitem{peract}
M.~Shridhar, L.~Manuelli, and D.~Fox,
``Perceiver-Actor: A multi-task transformer for robotic manipulation,''
in \emph{Proc. Conf. Robot Learning (CoRL)}, 2022.

\bibitem{octo}
Octo Model Team,
``Octo: An open-source generalist robot policy,''
in \emph{Proc. Robotics: Science and Systems (RSS)}, 2024.

\bibitem{saycan}
M.~Ahn \emph{et al.},
``Do as I can, not as I say: Grounding language in robotic
affordances,''
in \emph{Proc. Conf. Robot Learning (CoRL)}, 2022.

\bibitem{cap}
J.~Liang \emph{et al.},
``Code as policies: Language model programs for embodied control,''
in \emph{Proc. IEEE Int. Conf. Robotics and Automation (ICRA)}, 2023.

\bibitem{innermono}
W.~Huang \emph{et al.},
``Inner monologue: Embodied reasoning through planning with language
models,''
in \emph{Proc. Conf. Robot Learning (CoRL)}, 2022.

\bibitem{lora}
E.~J. Hu \emph{et al.},
``LoRA: Low-rank adaptation of large language models,''
in \emph{Proc. Int. Conf. Learning Representations (ICLR)}, 2022.

\bibitem{qlora}
T.~Dettmers, A.~Pagnoni, A.~Holtzman, and L.~Zettlemoyer,
``QLoRA: Efficient finetuning of quantized LLMs,''
in \emph{Proc. Advances in Neural Information Processing Systems
(NeurIPS)}, 2023.

\bibitem{soups}
M.~Wortsman \emph{et al.},
``Model soups: Averaging weights of multiple fine-tuned models improves
accuracy without increasing inference time,''
in \emph{Proc. Int. Conf. Machine Learning (ICML)}, 2022.

\bibitem{taskarith}
G.~Ilharco \emph{et al.},
``Editing models with task arithmetic,''
in \emph{Proc. Int. Conf. Learning Representations (ICLR)}, 2023.

\bibitem{ewc}
J.~Kirkpatrick \emph{et al.},
``Overcoming catastrophic forgetting in neural networks,''
\emph{Proc. National Academy of Sciences}, vol.~114, no.~13,
pp.~3521--3526, 2017.

\end{thebibliography}

\end{document}
